\section{Root Systems}

\begin{frame}{The Root System}
    From the Cartan-Weyl decomposition of a semi-simple Lie algebra $\mathfrak{g}$, the roots (eigenvalues of the adjoint action) form a highly symmetric geometric structure inside a Euclidean space $E = \mathfrak{h}_{\mathbb{R}}^\vee$.

    \vspace{0.3cm}
    \textbf{Physical Meaning:} The roots correspond to the non-Cartan generators of our symmetry (e.g., $W^\pm$ bosons in the Standard Model, acting as raising/lowering operators for quantum numbers).

    \vspace{0.5cm}
    The rigidity of quantum mechanics (specifically $SU(2)$ representation theory) forces these roots to satisfy strict geometric axioms.

    \vfill
    % Bottone ipertestuale verso l'appendice formale
    \hyperlink{backup_rootsystem}{\beamerbutton{Formal Definition of Root System}}
\end{frame}

\begin{frame}{The Crystallographic Restriction}
    The requirement that roots must fold into representations of independent $SU(2)$ subalgebras forces the quantity $n_{\alpha\beta} = \frac{2(\alpha, \beta)}{(\alpha, \alpha)}$ to be an integer.

    % OPTIONAL: Se il tempo stringe, enuncia solo la formula finale
    By the definition of the Euclidean scalar product, we can relate this to the angle $\theta_{\alpha\beta}$ between two roots $\alpha$ and $\beta$:
    \begin{equation}
        n_{\alpha\beta} n_{\beta\alpha} = 4 \frac{(\alpha, \beta)^2}{\|\alpha\|^2 \|\beta\|^2} = 4 \cos^2 \theta_{\alpha\beta} \in \mathbb{Z}
    \end{equation}

    Because $\cos^2 \theta \le 1$, the product must be strictly bounded:
    \begin{equation}
        n_{\alpha\beta} n_{\beta\alpha} \in \{0, 1, 2, 3, 4\}
    \end{equation}

    \vfill
    % Bottone ipertestuale verso l'appendice con la dimostrazione estesa
    \hyperlink{backup_proof}{\beamerbutton{Extended Algebraic Proof}}
\end{frame}

\begin{frame}{Allowed Angles and Lengths}
    Excluding the trivial collinear case ($4$), the allowed angles between any two linearly independent simple roots are heavily restricted.

    \vspace{0.3cm}
    For $\theta \ge \pi/2$, the only allowed geometric configurations are:
    \begin{itemize}
        \item $\theta = \pi/2$ ($90^\circ$): Roots are orthogonal.
        \item $\theta = 2\pi/3$ ($120^\circ$): Equal lengths ($\|\beta\|/\|\alpha\| = 1$).
        \item $\theta = 3\pi/4$ ($135^\circ$): Length ratio is $\sqrt{2}$.
        \item $\theta = 5\pi/6$ ($150^\circ$): Length ratio is $\sqrt{3}$.
    \end{itemize}

    \vspace{0.3cm}
    \textbf{Why it matters:} This mathematical rigidity is the profound reason why there are only a handful of distinct continuous gauge groups (like $SU(N)$, $SO(N)$) allowed in physics!
\end{frame}